% !TeX root = ../papers/nc-p3.tex
\section{分岐中心の各直線が掃く曲面は平面である}
\label{sec:plane-branch-pencils}
\begin{theorem}\label{thm:plane-branch-pencils}
標準正則窓で、正方向の直線ファイバーの中心が一本の直線 $L$ を含むとする。
その後続は一つの平面内の直線束であり、束の中心 $q$ での逆方向ファイバーは $L$ 全体である。
特にどの曲線族も二次曲面の母線を掃くことはない。
さらに $L$ の線型包を点座標に残し、関係六平面を制限すると三つの独立な関係が残る。
点座標で $[x_0:x_1]\in L$、後続平面を $b=0$、束の中心を $b=v=w=0$ とすれば
その関係は $x_0b,x_1b,x_0v+x_1w$ の包である。
\end{theorem}
\begin{proof}[滑らかな母線族の反対側の分離]
位置0を $L$、後続位置1を滑らかな二次曲面 $\mathbb P Y^*\times\mathbb P T^*$ とし
$V_1=Y\otimes T$ とする。前定理の最後の計算より、位置0を $L$ に射影した四テンソルは
$[x,y]f(t,z,u)$、$f$ の $u$ 支持次元は3である。$z,u$ は位置2,3を表す。
従って $R_1=Y\otimes K$ である。実際、$K\subset T\otimes V_2$ は $u$ 縮約による三空間で、
他の隣接切断の階数は $2\cdot3=6$ なので全六関係が得られる。

一般の $t\in\mathbb P T^*$ で $K(t)$ は階数3である。
さもなくば滑らかな二次曲面上で全ての $s\otimes t$ が順方向の直線ファイバーを持ってしまう。
よってその核点 $z(t)$ は $s$ に依存せず、逆方向に $\mathbb P Y^*\times\{t\}$ 全体を持つ。
全ファイバーは高々直線なので、異なる $t$ で同じ $z(t)$ は生じない。
その像の閉包は逆方向の分岐中心の直線 $L'$ である。その母線族は先と同じ二次曲面の他方の族となる。
同じ二項切断を逆向きに適用すれば、$R_0=J\otimes T$ である。
ここで逆向きの二項基底を改めて選んでも、二つの直線族による $V_1=Y\otimes T$ の分解は
各二空間の基底変更を除いて一定である。実際、滑らかな二次曲面の順序付き両直線族のパラメータが
それぞれの射影直線を決め、超平面束はその二つの一次束の外積である。

これで位置1の $Y,T$ の間で分離が起きる。前の関係は $V_0,Y$ だけに関わり、
次の関係は $T,V_2$ だけに関わる。四移動ごとに得られる完全な経路ブロック
$T,V_2,V_3,V_0,Y$ の次元を $b_0$ とすれば、その反復の次元は $b_0$ の冪と
次数に依らない両端ブロックの次元の積である。これは標準の三次Hilbert増大にならない。
\end{proof}

\begin{proof}[二次錐の母線族]
錐の頂点を $q$ とし、$\ker q\subset V_1$ を二項二次形式 $S^2(y)$ に同定する。
$L$ から母線のパラメータ直線への写像は同型である。
実際、単射かつ全射で、標数零の射影直線間のそのような射は次数1である。
位置0の二変数を $x$ として、制限関係 $R'_0$ は $[x,y]S^1(y)$ に等しい。
全ての関係は $q$ で消え、底二次曲線の $y=x$ で消えるのでこの二空間に入り、
各 $x$ で評価の階数が2なので等号である。
従って制限テンソルは $[x,y]f(y,z,u)$ となる。$f$ は $y$ の一次であり、
各 $x$ 評価では非零一次式の乗法だけが加わるので、$u$ 支持は一定である。
前定理の滑らかな場合の末尾と同じくその次元は3である。

$u$ で縮約した $f$ の三次元像を $B\subset S^1(y)\otimes V_2$ とする。
$S^1(y)B\subset R_1$ である（$x,u$ での全縮約を取る）。
これは $\ker q\otimes V_2$ に入るため六次元未満である。
もし六次元なら $R_1$ 全体がその部分に入り、位置1の全支持に反する。
従って $S^1\otimes B\to S^2\otimes V_2$ の核に非零元があり、
一次形式の互いに素な二元を $y_0,y_1$ として
$y_0 b_0+y_1 b_1=0$ となる。各係数で割ればある $0\ne d\in V_2$ に対して
$y_0d,y_1d\in B$ である。すると $S^2(y)\otimes kd\subset R_1$ は
同じ末尾文字の三つの独立な純関係で、定理\ref{thm:branching-centers-curve}の初段に反する。
\end{proof}

\begin{proof}[平面内の表示]
残る平面では、一般点の前点一意性により束となることは定理\ref{thm:branching-lines-only}で見た。
$L$ から束への写像も射影同型なので、後続平面上の入射は一つの非退化な二変数対の式で切れる。
$L$ への制限関係は従って表示した三空間に入る。二次元しかなければ、その二関係の
全点評価は階数2であるから、像束を $S$ として $2\mathcal O\simeq S(1)$ となる。
しかし $S$ は定数直線 $kb$ と母線パラメータの $\mathcal O(-1)$ の拡大であり、
$\det S(1)$ は $\mathcal O(1)$ である。よって三次元全部を得る。
\end{proof}

\begin{corollary}\label{cor:rank-three-plane-branch-slice}
上の平面の場合、制限四テンソルは独立な $C_0,C_1,C_2\in R_2$ により
\[
           b(x_0 C_0+x_1 C_1)+(x_0 v+x_1 w)C_2
\]
である。$C_2$ は純関係である。$D=L^\perp\subset V_0$ について
$\dim(R_0\cap(D\otimes V_1))=3$ であり、
$\dim(R_3\cap(V_3\otimes D))\ge2$ である。
特に純関係を持たない正則窓では、全射影の直線ファイバーの中心は有限個である。
\end{corollary}
\begin{proof}
$R'_0$ の次元から表示と三独立性、交叉次元3が従う。他の隣接切断は $x_0,x_1$ の両成分に
$C_2$ の $V_3$ 支持を含むので階数は $2\operatorname{rank} C_2$ 以上である。
一方、$V_1$ 側の支持は $\langle b,v,w\rangle$ に入り、六次元全部ではなく階数は5以下である。
これは $R_3$ を位置0で射影した階数でもある。
実際には $q$ の逆ファイバーが $L$ 全体なので $w(q)$ の位置0の支持は $D$ である。
位置2からの縮約は少なくとも二次元であり、$R_3\cap(V_3\otimes D)$ に入る。
従って射影階数は4以下である。$C_2$ の階数が2なら、その射影像は
$\langle x_0,x_1\rangle\otimes \operatorname{supp}_{V_3} C_2$ 全体に等しく、
異なる $L$ の点の逆端支持が同じ二平面となって不可能である。
\end{proof}
